% ============================================================
% §2.1 · Resumen Ejecutivo — Manual Técnico K-QGMIP
% ============================================================

\section{Resumen Ejecutivo}

\subsection{Problema abordado y su relevancia}

La \textbf{Teoría de la Información Integrada versión 4.0 (IIT 4.0)} propone que la consciencia
y la causalidad irreducible de un sistema pueden cuantificarse mediante el valor $\varphi$
(\textit{phi}). Para calcularlo es necesario encontrar la \textbf{Partición de Mínima Información
(MIP)}: la descomposición del sistema en subconjuntos que minimiza la pérdida de información
causal al tratar esas partes como independientes.

Las implementaciones previas del proyecto GeoMIP (estrategia geométrica) y QNodes (estrategia
submodular) resuelven este problema únicamente para \textbf{biparticiones} ($k = 2$), es decir,
para la división del sistema en exactamente dos partes. Sin embargo, la formulación completa de
IIT 4.0 exige explorar \textbf{k-particiones} ($k \geq 3$), lo cual introduce una explosión
combinatoria severa: el número de k-particiones distintas de $n$ variables está dado por el
\textbf{número de Stirling de segundo tipo} $S(n, k)$, que crece superexponencialmente.

\begin{table}[H]
  \centering
  \caption{Número de k-particiones posibles $S(n,k)$ para distintos valores de $n$ y $k$.
           Resaltados en gris los casos computacionalmente inviables por búsqueda exhaustiva.}
  \label{tab:stirling}
  \begin{tabular}{ccccc}
    \toprule
    $n$ & $k=2$ & $k=3$ & $k=4$ & $k=5$ \\
    \midrule
    5  & 15        & 25        & 10            & 1          \\
    8  & 127       & 966       & 2\,646        & 3\,025     \\
    10 & 511       & 9\,330    & 145\,750      & 1\,082\,250 \\
    20 & $\sim$524\,K & $\sim$580\,M & \cellcolor{gray!20}--- & \cellcolor{gray!20}--- \\
    25 & $\sim$16\,M  & \cellcolor{gray!20}--- & \cellcolor{gray!20}--- & \cellcolor{gray!20}--- \\
    \bottomrule
  \end{tabular}
\end{table}

\smallskip
El presente proyecto extiende ambas estrategias al caso general
$k \in \{2, 3, 4, 5\}$ bajo las denominaciones \textbf{KGeoMIP} y \textbf{KQNodes},
garantizando que para $k = 2$ los resultados son exactamente equivalentes a las implementaciones
originales.

\subsection{Enfoque algorítmico implementado}

Se implementaron dos estrategias complementarias de naturaleza distinta:

\subsubsection*{KGeoMIP — Refinamiento divisivo geométrico (Estrategia E4)}

KGeoMIP ancla la solución óptima de $k=2$ obtenida por GeoMIP (bipartición exacta via tabla de
costos BFS sobre el hipercubo de Hamming) y aplica un \textbf{refinamiento divisivo top-down}:

\begin{enumerate}
  \item Se construye una \textbf{matriz de similitud causal} $S \in \mathbb{R}^{D \times D}$
        a partir de la tabla de transiciones de GeoMIP, donde $S_{ij}$ cuantifica cuán
        similar es el comportamiento causal de las variables $i$ y $j$.
  \item En cada paso $t = 2, 3, \ldots, k$, de todas las partes actuales se selecciona
        la que produce el menor \textbf{incremento de pérdida} $\Delta\Phi$ al bipartirla,
        usando la similitud causal para guiar la selección de candidatos.
  \item El proceso se repite $k-1$ veces hasta obtener la k-partición deseada.
\end{enumerate}

Complejidad dominante: $\mathcal{O}(k \cdot D^2 \cdot 2^D)$, donde $D$ es la dimensión del
subsistema.

\subsubsection*{KQNodes — Refinamiento iterativo submodular (Criterio C4)}

KQNodes extiende el algoritmo de Queyranne (Ordenamiento de Máxima Adyacencia, MAO) al caso
$k \geq 3$ mediante un \textbf{refinamiento greedy con MinHeap}:

\begin{enumerate}
  \item Se construye un oracle restringido por bloque que aproxima el costo de bipartición
        mediante promedios ponderados sobre tensores n-dimensionales.
  \item En cada iteración del refinamiento, el MinHeap selecciona la parte con menor
        $\varphi_{\text{local}}$ (criterio C4) y la bipartite usando Queyranne (exacto dentro
        del bloque).
  \item La EMD-Effect real se evalúa una sola vez sobre la k-partición ganadora.
\end{enumerate}

Complejidad dominante: $\mathcal{O}(k \cdot D^3 \cdot N)$, donde $N$ es el número de nodos
y $D$ la dimensión del subsistema.

\subsection{Principales resultados y contribuciones}

\begin{itemize}
  \item \textbf{Regresión exacta $k=2$}: ambas estrategias producen resultados idénticos a
        GeoMIP y QNodes para $k=2$, verificado con tolerancia $|\Delta\varphi| < 10^{-9}$ sobre
        13 casos de prueba (redes $n \in \{5, 8, 10\}$).

  \item \textbf{Monotonicidad garantizada}: $\varphi(k+1) \geq \varphi(k)$ por construcción
        en ambas estrategias, propiedad consistente con la teoría IIT 4.0.

  \item \textbf{Escalabilidad}: KGeoMIP procesa sistemas de hasta $n = 25$ nodos y KQNodes
        de hasta $n = 20$ nodos en tiempos de ejecución del orden de segundos a minutos,
        frente a días o semanas que requeriría la búsqueda exhaustiva.

  \item \textbf{Criterio C4 vs C1 en KQNodes}: los experimentos A/B muestran que el criterio
        C4 (MinHeap por $\varphi_{\text{local}}$ mínimo) produce particiones con menor brecha
        respecto al óptimo que C1 (bipartir la parte más grande), confirmando la superioridad
        de la heurística basada en información integrada.
\end{itemize}

\subsection{Limitaciones encontradas y recomendaciones de uso}

\begin{itemize}
  \item \textbf{No optimalidad global para $k \geq 3$}: ninguna de las dos estrategias
        garantiza encontrar la k-partición de mínimo $\varphi$ global. Ambas son heurísticas
        de refinamiento cuya calidad depende de la calidad de la bipartición ancla ($k=2$) y
        del criterio de selección.

  \item \textbf{Crecimiento exponencial del oracle}: el oracle de KQNodes evalúa tensores
        de forma $(N, 2^m)$ donde $m$ es la dimensión del bloque; para $D \geq 22$ el consumo
        de memoria puede superar los 8 GB. Se recomienda usar KGeoMIP para sistemas grandes.

  \item \textbf{KGeoMIP limitado por la tabla BFS}: la tabla T de GeoMIP ocupa $\mathcal{O}(2^D)$
        en memoria. Para $D > 25$ se requiere memoización parcial o acceso a disco.

  \item \textbf{Recomendación de uso}: para $n \leq 15$ y $k \leq 5$ ambas estrategias son
        apropiadas; para $15 < n \leq 25$ se recomienda KGeoMIP-E4; para $n > 25$ se requieren
        extensiones fuera del alcance de este proyecto.
\end{itemize}
